% ====================================================================
% Apéndice A: Demostraciones sobre termomayorización
% ====================================================================
\chapter{Demostraciones sobre termomayorización}
\label{ap:thermomajorization_proofs}

\section{Demostración del teorema de Lu-Raz (\cref{thm:luraz_criterion})}

\begin{proof}
Sea $\vb{p}(t; \Tinit)$ la solución de la ecuación maestra con condición inicial 
$\pieq(\Tinit)$ y baño $\Tb$. Su expansión espectral es \cref{eq:luraz_expansion}:
\begin{equation}
    p_i(t; \Tinit) - \pieq_i(\Tb) = \sum_{k \geq 2} c_k(\Tinit) v_{i,k} e^{-\lambda_k t}.
\end{equation}
%
La distancia entrópica al equilibrio en el régimen asintótico está dominada por 
el modo de menor autovalor:
\begin{equation}
    \Dee(t; \Tinit) \approx \alpha \cdot c_2(\Tinit)^2 \, e^{-2 \lambda_2 t}, 
    \qquad \alpha > 0 \text{ constante.}
\end{equation}
%
Para dos preparaciones $\Th > \Tw$:
\begin{equation}
    \frac{\Dee(t; \Th)}{\Dee(t; \Tw)} \approx \frac{c_2(\Th)^2}{c_2(\Tw)^2}.
\end{equation}
%
\noindent Por tanto, $\Dee(t; \Th) < \Dee(t; \Tw)$ para $t$ grande si y solo si 
$\abs{c_2(\Th)} < \abs{c_2(\Tw)}$, que es \cref{eq:luraz_criterion}.
\end{proof}

\section{Demostración del teorema de Vu-Hayakawa (\cref{thm:vu_hayakawa})}

\begin{proof}[Esquema clásico]
Sean $\vb{p}, \vb{q}$ y $\gamma$ la distribución de Gibbs del baño. Necesitamos 
demostrar la equivalencia
\begin{equation}
    \vb{p} \succeq_{\Tb} \vb{q} \iff D_f(\vb{p} \| \gamma) \geq D_f(\vb{q} \| \gamma) 
    \quad \forall f \in \mathcal{F},
\end{equation}
%
donde $\mathcal{F}$ es la clase de funciones convexas con $f(1) = 0$.

\textbf{(}$\Rightarrow$\textbf{)}: Si $\vb{p} \succeq_{\Tb} \vb{q}$, entonces 
existe una matriz estocástica $\Phi$ con $\Phi \gamma = \gamma$ y 
$\Phi \vb{p} = \vb{q}$. Por la \emph{contractividad de Csiszar} de las 
$f$-divergencias bajo mapas que preservan la medida de referencia:
\begin{equation}
    D_f(\Phi \vb{p} \| \Phi \gamma) \leq D_f(\vb{p} \| \gamma) 
    \implies D_f(\vb{q} \| \gamma) \leq D_f(\vb{p} \| \gamma). \qed
\end{equation}

\textbf{(}$\Leftarrow$\textbf{)}: Esta dirección es no trivial y requiere la 
construcción explícita de la matriz $\Phi$. La idea de Lostaglio (2022) es 
emplear las \emph{Lorenz integrals}: definir
\begin{equation}
    \mathcal{I}_a(\vb{p}, \gamma) = \sum_i \pi(p_i; a, \gamma_i),
\end{equation}
%
para una familia uniparamétrica de funciones $\pi$ con propiedades adecuadas. 
Si $D_f(\vb{p} \| \gamma) \geq D_f(\vb{q} \| \gamma)$ para todo $f$, entonces en 
particular para la familia $\pi$, lo que implica que la curva de Lorenz de 
$\vb{p}$ está por encima de la de $\vb{q}$. Por el teorema clásico de 
Hardy-Littlewood-Pólya, existe la matriz biestocástica $\Phi$. \qed
\end{proof}

\section{Propiedad de la distancia entrópica de Lu-Raz}

\begin{proposition}
La distancia entrópica $\Dee$ del \cref{eq:dist_entropica} satisface
\begin{equation}
    \frac{d \Dee}{dt} \leq 0 \quad \text{en cualquier trayectoria markoviana},
\end{equation}
con igualdad solo en el equilibrio.
\end{proposition}

\begin{proof}
Diferenciando $\Dee = \sum_i p_i \ln(p_i / \pieq_i)$:
\begin{equation}
    \frac{d \Dee}{dt} = \sum_i \dot{p}_i \ln(p_i / \pieq_i) + \sum_i \dot{p}_i
    = \sum_{i,j} R_{ij} p_j \ln(p_i / \pieq_i),
\end{equation}
%
donde el segundo término se anula por conservación de la probabilidad. 
Reordenando con la condición de balance detallado $R_{ij} \pieq_j = R_{ji} \pieq_i$:
\begin{equation}
    \frac{d \Dee}{dt} = -\sum_{i < j} R_{ij} \pieq_j 
    \left[\frac{p_i}{\pieq_i} - \frac{p_j}{\pieq_j}\right]
    \ln\!\left[\frac{p_i/\pieq_i}{p_j/\pieq_j}\right].
\end{equation}
%
El integrando es siempre $\geq 0$ por $(x - y) \ln(x/y) \geq 0$. \qed
\end{proof}
